% Appendix. Fragment, \input by spacenav.tex inside \appendix.
\section{Why solvability cannot be checked by exhaustion}
\label{app:complexity}

Consider the most naive possible procedure for establishing that a level is
solvable, which is to enumerate every control sequence and simulate each
one. The control at each of $T$ steps is a direction and a thrust magnitude,
both continuous, so even with a discretization to $a$ distinct actions per
step the number of sequences is $a^T$, and at $T = 2250$ this is beyond
reach for any $a > 1$; the discrete nine-action set used here does not
rescue the situation ($9^{2250}$ has more than two thousand digits), and the
underlying continuous problem has no finite enumeration at all. Formal
results in motion planning make the same point: the generalized mover's
problem is PSPACE-hard \cite{reif1979,hopcroft1984,canny1988}, and remains
so with moving obstacles even for a point robot of bounded velocity
\cite{reif1994}; finite-horizon planning under partial observation is
PSPACE-complete and its infinite-horizon form undecidable
\cite{papadimitriou1987,madani2003}; and reachability for quite simple
classes of continuous and hybrid dynamics is undecidable
\cite{moore1990,henzinger1998,blondel2000}. Note that these are results
about problem classes, and say nothing about any single instance, and that
approximate search is not hopeless in principle, since kinodynamic planning
with bounded velocity and acceleration admits a polynomial-time
$\varepsilon$-approximation \cite{donald1993}; what they establish is that
no efficient general procedure is to be expected, and that nothing about
the field studied here, which couples every body's position to the craft's
acceleration at every step, places it among the restricted cases for which
an approximation is known.

Exhaustive checking is therefore replaced, in practice, by one of two
things, a simplifying assumption (a coarse model of the field, a restricted
family of control sequences, a short horizon) or a search that samples the
space and reports what it finds, and both are used in this paper. Both also
share a limitation that a chaotic regime makes acute. In a region where
nearby trajectories diverge exponentially, the numerical trajectory of a
given control sequence and the true trajectory of that sequence part company
within a few multiples of the Lyapunov time, since the shadowing of a
computed orbit by a true one breaks down in finite time when the local
expansion rate fluctuates \cite{hammel1987,sauer1997,sauer2002}, so that a
check which reports that a sequence arrives is a statement about a
trajectory that may correspond to nothing the craft can fly, and the precision of the check
becomes the limit before its cost does. The same literature supplies the
premise on which this paper's method rests, which is that ensemble
statistics over many orbits remain reliable where individual orbits do not
\cite{boekholt2015,portegieszwart2018}. Thus a level's difficulty is measured over a population of trajectories,
with the ensemble results above as the warrant, and the measurement
inherits whatever limits the probe has.
